\section{Introduction}
Streamflow forecasting plays a central role in water resources management, flood mitigation, and basin regulation and control\cite{GAO2026103148}. In recent years, deep learning techniques have achieved rapid progress in hydrological prediction\cite{milly2020colorado}, providing powerful predictive capabilities for streamflow forecasting\cite{apak2026multi,ZIPPER2026134909}. However, a critical question remains unresolved: do these models learn physically meaningful feature representations, or do they rely solely on complex statistical correlations to make predictions\cite{Núñez2023_XAI}? This distinction is of decisive importance for hydrological research. If models merely achieve statistical fitting, their generalization ability will be confined to the distribution of historical data, limiting their capacity to advance the understanding of the water cycle or support discoveries in Earth science. Conversely, if models are capable of encoding verifiable hydrological knowledge, deep learning can evolve from a mere fitting tool into a driver of scientific discovery\cite{MardShap2022}, enabling the identification of previously unresolved mechanisms in complex basins.

To address the “black-box” nature of deep learning hydrological models, constraint-based approaches represented by Physics-Informed Neural Networks (PINNs)\cite{SHI2025_PINN} attempt to embed hydrological governing equations into the loss function\cite{SCHMID2026124819}, thereby enforcing physical consistency in model outputs\cite{LIU2026_PINN}. Although this result-oriented integration of physical knowledge improves the plausibility of predictions to some extent, it does not fundamentally resolve the core issue of model interpretability\cite{Roscher2020_Explainable}. Even when model outputs conform to physical laws, the internal mechanisms by which input variables are utilized and features evolve remain hidden within a complex parameter space\cite{Krishnapriyan2021_PINN}. This implies that PINNs primarily mimic physical laws rather than actively discover hydrological processes\cite{NASER2026_Inrter}, and thus fail to reveal what hydrological mechanisms the model has truly learned. Consequently, their potential as tools for deep decision support and scientific knowledge discovery remains limited.

Another mainstream line of research investigates model interpretability through post-hoc analysis techniques, among which Local Interpretable Model-agnostic Explanations (LIME)\cite{Ribeiro2016_LIME_Origin} and SHapley Additive exPlanations (SHAP)\cite{lundberg2017_Shap_Origin} are the most widely used. These methods quantify the contributions of key hydrological variables—such as precipitation and evapotranspiration—to prediction outcomes based on cooperative game theory or local surrogate models\cite{Wagan2026_LIME}, thereby breaking the “black-box” barrier in a formal sense. However, such explanations are essentially post-hoc in nature. Their results heavily depend on the statistical properties of the data distribution rather than the model's intrinsic computational logic. More critically, they treat hydrological variables as isolated features, neglecting the prevalent collinearity among variables in natural systems. This mathematical decomposition, detached from the model's internal operations, often leads to unverifiable explanations and spurious correlations, making them unreliable for inferring hydrological scientific laws or understanding streamflow generation mechanisms in high-risk basins.

Overall, both PINNs and post-hoc interpretability techniques have made significant progress in improving the interpretability of streamflow forecasting models, offering feasible pathways to mitigate the “black-box” issue. Nevertheless, a fundamental scientific gap persists in the application of deep learning to hydrology: are the feature relationships extracted by models merely incidental statistical artifacts, or do they reflect robust hydrological physical laws? This gap directly hinders the paradigm shift of deep learning from a predictive tool to a scientific discovery instrument. Specifically, bridging this gap requires overcoming three interrelated systemic limitations in current interpretability research:

First, interpretability results lack objective verifiability\cite{NEURIPS2021_shap_residuals}, casting doubt on the validity of scientific discoveries derived from them. Unlike streamflow prediction results, which can be directly evaluated against observed ground truth using quantitative metrics, feature attribution-based interpretations cannot be directly validated and are often justified indirectly through high predictive accuracy\cite{TAKEFUJI2026_Shap_problem}. However, even highly accurate predictions may arise from complex but spurious statistical correlations\cite{e24050687_Shap_problem}, rather than physically meaningful hydrological knowledge. Without independent and quantitative validation of attribution results, it is difficult to determine whether models have truly learned physically meaningful relationships\cite{Benchmarkingofa_PhysicallyBasedHydrologicModel}. This severely limits the credibility of interpretability in high-stakes decision-making and impedes its role in scientific discovery.

Second, the lack of systematic evaluation of model behavior leads to a disconnect between interpretability results and actual hydrological physical logic. Current interpretability analyses often remain at the level of static feature importance rankings\cite{ZHOU2026103254,XIA2026103232}, neglecting the need to assess behavioral consistency under dynamic hydrological conditions\cite{10.2166/wcc.2017.149}. A scientifically meaningful hydrological model should respond to key drivers such as precipitation and evapotranspiration in accordance with fundamental hydrological principles\cite{hess-30-1077-2026}. For instance, a significant increase in precipitation should trigger a consistent rise in streamflow within the model\cite{nearing2021role}. Without such perturbation experiments or counterfactual validation\cite{simic2025perturbation,ZHANG2025123192}, interpretability results cannot be effectively linked to real hydrological behavior\cite{read2019process}, limiting their value for scientific understanding.

Finally, the lack of transferability validation across model architectures hinders the extraction of universal hydrological laws. The essence of scientific discovery lies in identifying generalizable principles that remain consistent across different observational tools or modeling approaches\cite{pmlr-v235-wang24j}. However, existing interpretability studies are often constrained by the parameterization of specific model architectures and rarely examine whether the identified hydrological relationships hold across different model families. If the extracted representations cannot be reproduced across models and scenarios, they are more likely to reflect overfitting to specific data distributions rather than true hydrological laws\cite{li2022transferability}, thereby limiting their capacity to support genuine scientific discovery.

Notably, in other scientific domains, systematic interpretability research paradigms have gradually matured, with standardized validation frameworks significantly enhancing the credibility of model explanations \cite{Haufe2026ExplainableAI}. However, such efforts remain relatively limited in hydrological modeling, and a unified and mature interpretability validation framework has yet to be established. The absence of such standardized validation paradigms has become the final critical barrier preventing deep learning-based hydrological models from evolving from predictive tools into engines of scientific discovery.

This study aims to transcend the descriptive limitations of traditional interpretability analysis by constructing a full-chain interpretability validation framework to systematically investigate the internal information propagation mechanisms of deep learning-based hydrological models. Rather than merely focusing on \textit{what the model predicts}, this work seeks to explore, at a fundamental level, \textit{how information flows within the model}, and \textit{how such internal knowledge supports model generalization under complex hydrological conditions}. To this end, the study focuses on the following four core scientific questions:

\begin{enumerate}

    \item \textbf{Structure-performance mapping mechanism:} What intrinsic relationships exist between model structure and predictive performance?

    \item \textbf{Existence and distribution of physical knowledge:} Do deep learning-based hydrological models encode physical knowledge, and if so, what are its spatial and feature-level distribution patterns?

    \item \textbf{Response consistency across hydrological scenarios:} Can the model generate physically consistent feature representations under different hydrological conditions?

    \item \textbf{Generality and transferability of scientific knowledge:} Is the physical knowledge learned by the model transferable across different scenarios and model architectures?

\end{enumerate}

To address these scientific questions, this paper proposes an interpretable hydrological model (STEH), aiming to achieve full-chain interpretability from input features, internal model structures, to output results. The main contribution of this study lies not only in the development of an interpretable deep learning framework, but also the establishment of a reference path for hydrological model interpretability verification. Framed as \textbf{\textit{A Deep Dive into Model Structure}}, this work integrates structural analysis, behavioral verification, and cross-model evaluation to move interpretability beyond descriptive attribution toward rigorous assessment of whether model knowledge is \textbf{\textit{verifiable}}, physically meaningful, and \textbf{\textit{transferable}}. In doing so, it provides preliminary insights into a fundamental question: \textbf{\emph{what kind of hydrological knowledge is learned, and is it verifiable and transferable?}}